\documentclass[12pt, a4paper]{article}

\usepackage[top=2.5cm, bottom=2.5cm, left=3cm, right=3cm]{geometry}
\usepackage{mdframed}
\usepackage[utf8]{inputenc}
\usepackage[T1]{fontenc}
\usepackage[spanish]{babel}
\usepackage{amsmath}
\usepackage{float}
\usepackage[margin=10pt, font=small, labelfont=bf]{caption}
\usepackage{booktabs}
\usepackage{graphicx}
\usepackage{xcolor}
\usepackage{hyperref}
\usepackage{natbib}
\usepackage{longtable}
\usepackage{comment}
\usepackage{setspace}
\usepackage{titlesec}
\usepackage{fancyhdr}

\setlength{\textfloatsep}{18pt plus 2pt minus 4pt}
\setlength{\floatsep}{12pt plus 2pt minus 2pt}
\setlength{\intextsep}{14pt plus 2pt minus 2pt}
\setstretch{1.15}
\setlength{\parindent}{0pt}
\setlength{\parskip}{8pt}

\titleformat{\section}
  {\large\bfseries}{\thesection.}{0.5em}{}
  [\vspace{2pt}\hrule\vspace{4pt}]
\titleformat{\subsection}
  {\normalsize\bfseries}{\thesubsection.}{0.5em}{}
\titlespacing{\section}{0pt}{18pt}{8pt}
\titlespacing{\subsection}{0pt}{14pt}{6pt}

\hypersetup{colorlinks=true, linkcolor=blue!60!black}

\pagestyle{fancy}
\fancyhf{}
\fancyhead[L]{\footnotesize Centro Javeriano de Competitividad}
\fancyhead[R]{\footnotesize \reportnumber}
\cfoot{\thepage}
\setlength{\headheight}{14pt}
\renewcommand{\headrulewidth}{0.4pt}

\geometry{margin=1in}
\graphicspath{{Paper/figures/}{figures/}}
\newcommand{\tablepath}{Paper/tables/}
\newcommand{\figurelang}{es}
\newcommand{\reportseries}{Informes sobre la Productividad en Colombia}
\newcommand{\reportseriesdesc}{Evidencia para entender y fomentar la productividad}
\newcommand{\reportnumber}{Informe 01}
\newcommand{\reporttitle}{Dinámicas de la Productividad Laboral por Actividad Económica en Colombia 2010--2025}
\newcommand{\reportsubtitle}{}
\newcommand{\reportauthorone}{Oliver Pardo}
\newcommand{\reportauthoroneaffil}{Profesor asociado del Departamento de Administración, Pontificia Universidad Javeriana y Director del Centro Javeriano de Competitividad}
\newcommand{\reportauthortwo}{Jorge Orozco}
\newcommand{\reportauthortwoaffil}{Consultor de investigación del Centro Javeriano de Competitividad, Pontificia Universidad Javeriana}
\renewcommand{\contentsname}{Tabla de contenido}
\newif\ifpolicyreport
\policyreporttrue

\title{\reporttitle\\\reportsubtitle}
\author{\reportauthorone \and \reportauthortwo}
\date{\today}

\begin{document}

\begin{titlepage}
\thispagestyle{empty}
\vspace*{0.6cm}

\begin{center}

\includegraphics[width=0.58\textwidth]{logo_cjc_palatino_horizontal.png}
\end{center}

\vspace{1.5cm}

\begin{center}
{\Large\bfseries \reportseries\par}
\vspace{0.15cm}
{\large \reportseriesdesc\par}

\vspace{1.4cm}

{\large\bfseries \reportnumber\par}
\vspace{0.45cm}
{\Huge\bfseries \reporttitle\par}
\vspace{0.35cm}
{\Large \reportsubtitle\par}
\vspace{0.65cm}
{\large\bfseries \reportauthorone\par}
{\normalsize \reportauthoroneaffil\par}
\vspace{0.25cm}
{\large\bfseries \reportauthortwo\par}
{\normalsize \reportauthortwoaffil\par}
\end{center}

% El recuadro de propósito de la serie se omite en la portada para mantenerla
% en una sola página después de incluir autores y afiliaciones.

\vfill

\begin{center}
%{\large Centro Javeriano de Competitividad\par}
{\large \today\par}
\end{center}

\end{titlepage}

\begin{mdframed}[
  linewidth=1pt,
  linecolor=gray!60,
  backgroundcolor=gray!6,
  innertopmargin=12pt,
  innerbottommargin=12pt,
  innerleftmargin=14pt,
  innerrightmargin=14pt
]
\section*{Resumen ejecutivo}
\addcontentsline{toc}{section}{Resumen ejecutivo}
\textbf{Este informe analiza la evolución de la productividad laboral en Colombia entre 2010 y 2025.} Para ello combina información del Producto Interno Bruto (PIB) con información de ocupados y horas trabajadas de la Gran Encuesta Integrada de Hogares (GEIH). El objetivo es medir cuánto produce la economía por cada trabajador y por cada hora trabajada, cuánto ha crecido esa productividad durante los últimos quince años y cómo ha variado entre las ramas de actividad económica.

\textbf{El crecimiento del PIB se descompone entre el crecimiento de la población ocupada, el crecimiento de las horas trabajadas por cada persona ocupada y el crecimiento del producto durante cada hora trabajada.} El PIB real creció a una tasa anual de más del 3\% entre 2010 y 2025, explicada por un aumento en el número de ocupados de más del 2\% anual, una disminución de las horas por trabajador de casi 0,5\% anual y un aumento del PIB por hora trabajada del 1,5\% anual. Como resultado, el PIB por trabajador ha venido creciendo al 1\% anual, más despacio que el PIB por hora trabajada.

\textbf{La productividad por actividad económica muestra dinámicas muy heterogéneas.} Actividades como información y comunicaciones han mostrado crecimientos del 4,5\% anual en el PIB por hora trabajada durante los últimos quince años, mientras que otras como la minería han mostrado caídas del 3,1\% anual. Las razones para este comportamiento tan diverso escapan del alcance de este informe, pero serán objeto de análisis en próximos informes del Centro Javeriano de Competitividad.

\textbf{El crecimiento de la ocupación tendió a concentrarse en actividades donde la productividad creció menos.} Esto no es necesariamente una señal de una mala asignación del empleo, dado que algunas de las actividades donde la productividad creció más contaban al inicio del periodo con una productividad relativamente baja. No obstante, el hecho de que las actividades donde más creció la ocupación sean precisamente las actividades donde menos creció la productividad señala una oportunidad perdida para acelerar el crecimiento económico.
\end{mdframed}

\newpage
\tableofcontents
\newpage

\section{Introducción}

\textbf{La productividad ocupa un lugar central en cualquier agenda de desarrollo.} La literatura económica sobre crecimiento ha mostrado que, en el largo plazo, la capacidad de sostener aumentos en los ingresos de los habitantes depende menos de acumular más capital o de incorporar más trabajadores y horas trabajadas, y más de producir mayor valor agregado con los recursos disponibles. En otras palabras, el crecimiento económico sostenido requiere aumentos de productividad.

\textbf{El objetivo de este informe es medir la evolución de la productividad laboral en Colombia entre 2010 y 2025.} Con este fin se combina el PIB total y el PIB por actividad económica reportado por el DANE con la información de ocupados y horas trabajadas de la Gran Encuesta Integrada de Hogares (GEIH). A partir de esa combinación se construyen dos indicadores complementarios: el \textbf{PIB por trabajador} y el \textbf{PIB por hora trabajada}, tanto para la economía agregada como para las actividades económicas. La distinción es importante porque el número de ocupados y las horas promedio no necesariamente crecen al mismo ritmo, lo cual hace que las conclusiones que se derivan de un indicador u otro no sean necesariamente las mismas.

\textbf{En Colombia, el PIB por hora trabajada ha crecido más que el PIB por trabajador}. Durante los últimos quince años la economía logró crecer a una tasa promedio de más del 3\% anual. Sin embargo, es importante desagregar los motivos de esa expansión. Mientras la población ocupada creció a más del 2\% anual, el PIB por trabajador aumentó apenas cerca de 1\% anual. Al mismo tiempo, las horas semanales promedio por trabajador disminuyeron casi 0,5\% anual. Esa reducción hace que el PIB por hora haya crecido más que el PIB por trabajador. En efecto, el crecimiento del PIB por hora fue de 1,5\% anual.

\textbf{Existe una gran dispersión en el crecimiento de la productividad entre agrupaciones de actividad económica.} Este informe no solo estudia la dinámica de la productividad agregada, sino también la productividad de trece ramas comparables de actividad económica. Durante el periodo 2010--2025, actividades como arte, entretenimiento, recreación y otros servicios alcanzaron un crecimiento anualizado superior al 7\% en el PIB por hora trabajada, mientras que actividades como explotación de minas y canteras enfrentaron una reducción anualizada del 3,1\% en su PIB por hora trabajada.

\textbf{Una pregunta que se aborda es si la ocupación creció en las actividades económicas donde también creció la productividad.} Los resultados muestran que, entre 2010 y 2025, las actividades con el mayor aumento porcentual en el número de ocupados no fueron aquellas donde más creció el PIB por trabajador o el PIB por hora trabajada. Esto no necesariamente debe interpretarse como una señal de mala asignación de los ocupados, dado que algunas de las actividades con mayor crecimiento de la productividad partían de niveles relativamente bajos. Por esa razón, informes posteriores del Centro Javeriano de Competitividad estudiarán con más detalle la asignación de los factores de producción entre distintas actividades económicas, con el fin de estimar las potenciales ganancias de eficiencia que se alcanzarían al reducir las barreras que impiden que los recursos se asignen hacia sus usos más productivos.

\textbf{El informe está organizado en seis secciones, incluyendo la presente introducción.} La sección 2 presenta la metodología y la descomposición que permite separar el crecimiento del PIB entre ocupación, horas trabajadas y productividad por hora. La sección 3 describe la evolución agregada de la productividad laboral en Colombia. La sección 4 analiza la evolución de la productividad de cada una de las trece ramas comparables de actividad económica. La sección 5 estudia la relación entre crecimiento de la ocupación y crecimiento de la productividad a través de actividades económicas. Finalmente, la sección 6 presenta las conclusiones. El apéndice incluye la descomposición del crecimiento por actividad económica y la evolución de la productividad para una desagregación aún mayor.

\begin{comment}
\subsection{Comparación internacional: productividad laboral}

Antes de analizar la evolución del empleo y de los ingresos laborales en Colombia, conviene ubicar el desempeño del país en perspectiva internacional. Una forma sintética de hacerlo es comparar el PIB por trabajador, una medida amplia de productividad laboral que aproxima el valor agregado generado por cada persona ocupada. Aunque este indicador no captura todas las diferencias en calidad del empleo, informalidad, horas trabajadas o composición por actividad económica, permite establecer si el crecimiento colombiano ha sido suficiente para cerrar brechas frente a economías comparables y frente a referentes de mayor ingreso.

\begin{table}[H]
\centering
\caption{PIB por trabajador y crecimiento anualizado, 2010--2024}
\label{tab:pib_trabajador_internacional}
\begin{tabular}{lrrr}
\toprule
País & 2010 & 2024 & Crec. anualizado \\
\midrule
Costa Rica & 45.175 & 61.875 & 2,27\% \\
República Dominicana & 40.301 & 53.911 & 2,10\% \\
Perú & 22.321 & 29.844 & 2,10\% \\
Colombia & 32.469 & 40.363 & 1,57\% \\
Estados Unidos & 129.284 & 153.544 & 1,24\% \\
Chile & 55.677 & 63.924 & 0,99\% \\
Brasil & 39.603 & 41.441 & 0,32\% \\
México & 49.943 & 48.671 & -0,18\% \\
\bottomrule
\end{tabular}
\caption*{\footnotesize Nota: valores en dólares internacionales constantes de 2021 ajustados por paridad de poder adquisitivo (PPP). El crecimiento anualizado corresponde a la variación compuesta entre 2010 y 2024. Fuente: Banco Mundial, World Development Indicators, indicador \texttt{SL.GDP.PCAP.EM.KD}.}
\end{table}

\textbf{En perspectiva internacional, Colombia muestra un avance real de la productividad laboral, pero insuficiente para cerrar de manera clara sus brechas frente a varios pares regionales.} Entre 2010 y 2024, el PIB por trabajador en Colombia pasó de 32.469 a 40.363 dólares internacionales constantes de 2021, equivalente a un crecimiento anualizado de 1,57\%. Esta tasa supera la observada en Estados Unidos, Chile, Brasil y México durante el mismo período, un resultado que a primera vista puede parecer sorprendente. Sin embargo, debe interpretarse a la luz de los niveles iniciales: Estados Unidos y Chile partían de un PIB por trabajador sustancialmente más alto que Colombia, por lo que un crecimiento relativo menor es consistente con la idea, ampliamente documentada en la literatura económica, de que las economías con menor productividad inicial pueden crecer más rápido cuando tienen espacio para adoptar tecnologías, reorganizar procesos productivos y reasignar recursos hacia actividades de mayor valor agregado.

La comparación también muestra los límites de esa lectura optimista. Costa Rica y República Dominicana ya registraban en 2010 un PIB por trabajador superior al colombiano y, aun así, crecieron más rápido entre 2010 y 2024. Esto sugiere que el rezago de Colombia no se explica únicamente por su punto de partida, sino por la velocidad con la que transforma mejoras en empleo, capital humano, formalización y estructura productiva en aumentos sostenidos de productividad. En otras palabras, Colombia ha avanzado, pero no al ritmo necesario para converger con las economías latinoamericanas que han logrado combinar mayores niveles iniciales con una dinámica reciente más vigorosa.
\end{comment}


\section{Metodología}

\textbf{La medición de productividad laboral combina información de cuentas nacionales y de la GEIH.} El PIB anual se toma de las cuentas nacionales del DANE por el enfoque de producción y se expresa en pesos constantes de 2015. El número de ocupados se calcula como el promedio anual de personas ocupadas mensualmente de acuerdo a la GEIH y se valida con los totales oficiales de ocupación del DANE.

\textbf{Las horas trabajadas también se construyen a partir de la GEIH.} Primero se calcula el total de horas semanales trabajadas por los ocupados, usando los factores de expansión de la encuesta. Luego ese total semanal se multiplica por 52 semanas para obtener una medida anual de horas trabajadas. Así, el PIB por trabajador divide el PIB anual por el promedio anual de ocupados, mientras que el PIB por hora trabajada divide el PIB anual por el total anual de horas trabajadas.

\textbf{Para el análisis por actividad económica, el PIB del DANE y la información laboral de la GEIH deben estar expresados en agrupaciones comparables.} La medición principal se hace con \textbf{13 ramas comparables} de la Clasificación Industrial Internacional Uniforme de todas las actividades económicas, Revisión 4 Adaptada para Colombia (CIIU Rev. 4 A.C.). %Estas ramas parten de las agrupaciones de cuentas nacionales del DANE, pero separan la agrupación R+S+T en dos: arte, entretenimiento, recreación y otros servicios, por un lado, y hogares como empleadores, por el otro. Esta apertura evita mezclar dinámicas económicas muy distintas y es posible porque tanto las cuentas nacionales como la GEIH permiten construir una medición comparable para ambas ramas. 
Como una parte de la GEIH del periodo analizado fue clasificada con CIIU Rev. 3 A.C., se usa la tabla correlativa oficial entre CIIU Rev. 3 A.C. y CIIU Rev. 4 A.C. para homologar ocupados y horas a las agrupaciones usadas por las cuentas nacionales. %Cuando la correlativa identifica una reasignación clara de una clase Rev. 3 a una agrupación Rev. 4, los ocupados y las horas se asignan a esta última. Cuando la correspondencia es parcial y una misma clase Rev. 3 se reparte entre varias actividades Rev. 4, no se distribuyen ocupados ni horas de manera mecánica.

\textbf{A lo largo del informe se usan las expresiones PIB por actividad económica, PIB por trabajador y PIB por hora.} En sentido estricto, cuando el análisis se desagrega por actividad económica, el numerador del respectivo indicador de productividad corresponde al valor agregado bruto de cada actividad. La diferencia entre el PIB total y el valor agregado bruto por actividad proviene de los impuestos netos sobre los productos, los cuales no están desagregados por actividad económica. Sin embargo, se usan las expresiones PIB por trabajador y PIB por hora por conveniencia pedagógica, aún cuando en realidad representan el valor agregado bruto por trabajador y el valor agregado bruto por hora trabajada. 

%\textbf{El DANE también publica esta información con mayor desagregación, a 25 y 61 agrupaciones.} Para usar esas aperturas, se busca el nivel más detallado en el que el PIB de cuentas nacionales y la información laboral de la GEIH pueden compararse sin repartir artificialmente ocupados u horas. Cuando el código de actividad económica a cuatro dígitos de la GEIH permite una homologación clara con una subactividad del DANE, los ocupados y las horas se asignan a esa subactividad; esto permite separar, por ejemplo, café, ganadería, carbón, petróleo, edificaciones, obras civiles y algunas actividades de transporte. Cuando esa correspondencia no es suficientemente limpia, las subactividades del DANE se suman y se presentan como una sola observación comparable. Por eso, el apéndice construido a partir de las 61 agrupaciones del DANE contiene menos de 61 observaciones: no se calculan productividades separadas cuando no existe una medición laboral separada y confiable.


\textbf{El crecimiento del PIB puede desagregarse entre el crecimiento de la ocupación, el crecimiento de las horas por trabajador y el crecimiento de la productividad por hora.} En primer lugar, el PIB puede escribirse como el producto entre el PIB por trabajador y el número de ocupados:
\begin{equation}\label{eq:pibtotal}
\text{PIB} = \text{Ocupados}\times
\left(\frac{\text{PIB}}{\text{Ocupados}}\right).    
\end{equation}
A su vez, el PIB por trabajador combina dos elementos: cuánto se produce por cada hora trabajada y cuántas horas trabaja, en promedio, cada persona ocupada:
\begin{equation}\label{eq:pibxocu}
\frac{\text{PIB}}{\text{Ocupados}} =
\left(\frac{\text{Horas}}{\text{Ocupados}}\right)
\times
\left(\frac{\text{PIB}}{\text{Horas}}\right).    
\end{equation}
Por lo tanto, el crecimiento del PIB por persona ocupada es aproximadamente la suma entre el crecimiento de las horas trabajadas por persona ocupada y el crecimiento del PIB por hora trabajada. De las ecuaciones (\ref{eq:pibtotal}) y (\ref{eq:pibxocu}) se deriva que el PIB puede descomponerse en tres componentes:
\begin{equation}
\text{PIB} = \text{Ocupados} \times \left(\frac{\text{Horas}}{\text{Ocupados}}\right) \times \left(\frac{\text{PIB}}{\text{Horas}}\right).    
\end{equation}
%En términos de tasas de crecimiento anualizadas, esta identidad implica:
%\[
%1 + g_{\text{PIB}} =
%\left(1 + g_{\text{PIB/Horas}}\right)
%\left(1 + g_{\text{Horas/Trabajador}}\right)
%\left(1 + g_{\text{Ocupados}}\right).
%\]
%Para tasas moderadas, la misma relación puede leerse de manera aproximada como una suma:
%\[
%g_{\text{PIB}} \approx g_{\text{PIB/Horas}} + g_{\text{Horas/Trabajador}} + g_{\text{Ocupados}}.
%\]
Así, el crecimiento del producto puede aproximarse por la suma del crecimiento en el número de ocupados, el crecimiento en las horas por trabajador y el crecimiento del producto por hora trabajada. %Esta descomposición permite distinguir si el crecimiento económico refleja principalmente una expansión del empleo o mejoras en la capacidad productiva del trabajo.

\textbf{Los resultados, los insumos y el código de replicación están disponibles en el repositorio del proyecto: \url{https://github.com/Kin-George/CJC-Monitor}.}

\section{Productividad laboral agregada}

\textbf{Esta sección presenta una mirada agregada a la productividad laboral en Colombia.} El objetivo es mostrar cómo evolucionaron, entre 2010 y 2025, el PIB real, el número de ocupados, las horas trabajadas y dos medidas complementarias de productividad: el PIB por trabajador y el PIB por hora trabajada. 

\begin{table}[H]
\centering
\caption{PIB, ocupados, horas y productividad laboral, 2010--2025}
\label{tab:ocupados_horas_resumen}
\small
\begin{tabular}{lrrr}
\toprule
Indicador & 2010 & 2025 & Crec. anual \\
\midrule
PIB real (Billones de pesos de 2015) & 640 & 1.021 & 3,16\% \\
Ocupados (Millones) & 16,8 & 23,0 & 2,09\% \\
PIB por trabajador (Millones de pesos de 2015) & 38,0 & 44,5 & 1,05\% \\
Horas semanales por trabajador  & 46,8 & 43,7 & -0,45\% \\
PIB por hora trabajada (Miles de pesos de 2015) & 15,6 & 19,6 & 1,51\% \\
\bottomrule
\end{tabular}
\caption*{\footnotesize Fuente: cálculos propios con base en Cuentas Nacionales y GEIH del DANE.}
\end{table}


\begin{figure}[H]
  \centering
  \includegraphics[width=0.92\textwidth]{Paper/figures/fig_pib_geih_descomposicion_crecimiento.png}
  \caption{Descomposición del crecimiento anual del PIB, 2010--2025}
  \label{fig:pib_geih_descomposicion_crecimiento}
  \caption*{\footnotesize Nota: las barras intermedias muestran contribuciones aproximadas en puntos porcentuales al crecimiento anualizado del PIB real. La descomposición parte de la identidad \(\text{PIB} = \text{ocupados} \times \text{horas por trabajador} \times \text{PIB por hora trabajada}\). Por redondeo, las contribuciones pueden no sumar exactamente al crecimiento reportado. Fuente: cálculos propios con base en Cuentas Nacionales y GEIH del DANE.}
\end{figure}

\textbf{El Cuadro \ref{tab:ocupados_horas_resumen} presenta, tanto para el inicio como para el final del periodo 2010--2025, el PIB real, el número de ocupados, el PIB por trabajador, las horas semanales por trabajador y el PIB por hora trabajada; la Figura \ref{fig:pib_geih_descomposicion_crecimiento} resume la contribución de cada componente al crecimiento anual del PIB real.} Primero, el PIB por trabajador se obtiene dividiendo el PIB real entre el número de ocupados: en 2025, una economía con una producción anual de 1.021 billones de pesos de 2015 y 23,0 millones de ocupados produce 44,5 millones de pesos anuales por trabajador. Segundo, el PIB por hora se obtiene dividiendo el PIB entre el total anual de horas trabajadas. Tercero, el PIB por trabajador también puede leerse como el producto entre el PIB por hora y las horas trabajadas por ocupado. Así, el crecimiento anual del PIB por trabajador combina un aumento del PIB por hora de 1,51\% con una caída de las horas semanales por trabajador de 0,45\%, lo que resulta en un crecimiento de 1,05\% anual. Finalmente, el crecimiento del PIB total combina ese 1,05\% anual de PIB por trabajador con el crecimiento de los ocupados, de 2,09\% anual, para llegar al 3,16\% anual del PIB real.

\textbf{En conclusión, del crecimiento anual de más del 3\% del PIB en los últimos quince años, aproximadamente 2 puntos porcentuales se explican por el aumento de la población ocupada, -0,5 puntos porcentuales se explican por una caída en las horas trabajadas por ocupado, y 1,5 puntos porcentuales se explican por un aumento en la productividad por hora.} El aumento de la productividad por hora es notable, especialmente si se compara con el modesto crecimiento de la productividad por trabajador del 1\% anual. Esta brecha entre el crecimiento de la productividad por hora y el crecimiento de la productividad por trabajador se explica por la caída no despreciable de las horas trabajadas por ocupado. En otras palabras, el crecimiento de la productividad por trabajador enmascara parcialmente una mejora aún mayor en el producto generado por cada hora trabajada.

\section{Productividad laboral por grandes ramas de actividad económica}

La productividad agregada oculta diferencias sustanciales entre actividades económicas. Para examinarlas, se cruza el PIB por actividad económica del DANE con los ocupados y horas de la GEIH asignados a ramas comparables de actividad económica. Esta aproximación permite estimar, para cada actividad económica, el PIB por trabajador y el PIB por hora trabajada entre 2010 y 2025, como se reporta en el Cuadro \ref{tab:pib_geih_productividad_sector}.

\begin{table}[H]
\centering
\caption{Productividad laboral por actividad económica CIIU, 2010--2025}
\label{tab:pib_geih_productividad_sector}
\footnotesize
\begin{tabular}{lrrrrrr}
\toprule
& \multicolumn{3}{c}{PIB por trabajador} & \multicolumn{3}{c}{PIB por hora} \\
& \multicolumn{3}{c}{\footnotesize Millones de pesos de 2015} & \multicolumn{3}{c}{\footnotesize Miles de pesos de 2015} \\
\cmidrule(lr){2-4}\cmidrule(lr){5-7}
Actividad económica & 2010 & 2025 & Crec. & 2010 & 2025 & Crec. \\
\midrule
Arte, recreación y otros servicios & 11,4 & 32,8 & 7,32\% & 5,6 & 16,0 & 7,28\% \\
Información y comunicaciones & 45,4 & 78,6 & 3,73\% & 17,6 & 34,0 & 4,49\% \\
Agropecuario & 13,5 & 20,4 & 2,79\% & 5,8 & 9,4 & 3,27\% \\
Hogares como empleadores & 6,4 & 9,4 & 2,63\% & 2,7 & 4,5 & 3,49\% \\
Adm. pública, educación y salud & 43,9 & 58,4 & 1,92\% & 19,5 & 26,1 & 1,95\% \\
Financieras & 101,0 & 124,5 & 1,40\% & 43,0 & 54,2 & 1,56\% \\
Comercio, transporte y alojamiento & 20,9 & 24,1 & 0,97\% & 7,9 & 10,1 & 1,64\% \\
Manufactura & 40,1 & 45,0 & 0,77\% & 16,5 & 19,5 & 1,12\% \\
Inmobiliarias & 280,0 & 284,6 & 0,11\% & 105,3 & 119,5 & 0,85\% \\
Profesionales y administrativas & 54,3 & 42,4 & -1,63\% & 24,8 & 20,6 & -1,22\% \\
Construcción & 41,8 & 26,1 & -3,08\% & 16,4 & 11,0 & -2,63\% \\
Minas & 200,8 & 115,8 & -3,60\% & 77,6 & 48,4 & -3,09\% \\
Servicios públicos & 176,5 & 99,5 & -3,75\% & 70,0 & 43,5 & -3,13\% \\
\bottomrule
\end{tabular}
\caption*{\footnotesize Nota: El numerador de ambos indicadores de productividad corresponde estrictamente al valor agregado bruto de cada actividad. %Las actividades económicas parten de las agrupaciones CIIU Rev. 4 A.C. del DANE; la agrupación R+S+T se separa en R+S y T porque las cuentas nacionales y la GEIH permiten construir una medida laboral comparable para ambas. 
Los ocupados y las horas se agregan desde GEIH usando la actividad económica reportada en la encuesta. Fuente: cálculos propios con DANE y GEIH.}
\end{table}

\begin{figure}[H]
  \centering
  \makebox[\textwidth][c]{\includegraphics[width=1.08\textwidth]{Paper/figures/fig_pib_geih_productividad_burbujas_2025.png}}
  \caption{Nivel, crecimiento y tamaño económico de la productividad por hora en 13 grandes ramas de actividad económica}
  \label{fig:pib_geih_productividad_burbujas_2025}
\end{figure}

\textbf{La Figura \ref{fig:pib_geih_productividad_burbujas_2025} ordena las 13 grandes ramas comparables de actividad económica según su nivel de productividad por hora en 2025 y su crecimiento entre 2010 y 2025.} El eje horizontal mide el PIB por hora trabajada en 2025 y el eje vertical mide su crecimiento anual. La línea azul vertical marca el nivel agregado de productividad por hora en 2025 y separa las ramas con niveles superiores e inferiores al agregado. La línea azul horizontal marca el crecimiento agregado de la productividad por hora entre 2010 y 2025 y separa las ramas que crecieron por encima y por debajo del agregado. Así, las ramas \emph{líderes en auge} combinan niveles y crecimientos superiores al agregado; el caso más claro es información y comunicaciones. Las \emph{aceleradoras} crecen por encima del agregado aunque todavía tienen niveles inferiores de productividad; allí sobresale arte, recreación y otros servicios. Las \emph{líderes en declive} conservan niveles altos, pero crecen por debajo del agregado; minas y servicios públicos son los casos más llamativos porque no solo crecieron por debajo del agregado, sino que redujeron su productividad por hora. Finalmente, las \emph{rezagadas} combinan niveles y crecimientos inferiores, como construcción.

%La Figura \ref{fig:pib_geih_descomposicion_sector}, incluida en el apéndice, descompone el crecimiento anual del PIB de cada actividad económica en los mismos tres componentes del ejercicio agregado: ocupados, horas por trabajador y PIB por hora trabajada. La comparación muestra que las horas por trabajador restan crecimiento en la mayoría de actividades, mientras que el peso relativo de la ocupación y de la productividad por hora varía sustancialmente entre actividades económicas.

\textbf{La descomposición por actividad económica revela una gran diversidad en el crecimiento de la productividad.} Las mayores tasas de crecimiento anual del PIB por trabajador se observan en arte, entretenimiento, recreación y otros servicios (7,32\% anual), información y comunicaciones (3,73\%) y agropecuario (2,79\%). Paralelamente se observa una caída en el producto por trabajador en actividades como servicios públicos (-3,75\% anual), minería (-3,60\% anual) y construcción (-3,08\% anual). En todos los casos, el PIB por hora trabajada evoluciona de manera más favorable que el PIB por trabajador, de donde se infiere una caída de las horas trabajadas por trabajador en todas las grandes ramas de actividad económica.

\textbf{La trayectoria anual de la productividad muestra distintos momentos para cada una de las grandes ramas de actividad económica.} Como se observa en la Figura \ref{fig:pib_geih_productividad_sector_indices} del apéndice, la productividad de minería aumentó hasta 2016 y luego cayó de manera persistente, en un patrón que coincide temporalmente con el final del ciclo de precios altos de las materias primas. En manufactura, la productividad aumentó hasta 2021 y luego retrocedió de forma continua hasta 2025. En información y comunicaciones, el crecimiento fue acelerado hasta 2016, momento a partir del cual se desaceleró. Por último, en arte, entretenimiento, recreación y otros servicios, la mayor parte del aumento se observa después de 2019. Una hipótesis para explorar en informes posteriores es si parte de este dinamismo está asociado con actividades de juegos de azar y apuestas en línea. Las razones específicas de estas distintas dinámicas escapan del alcance del presente informe, pero serán objeto de informes posteriores por parte del Centro Javeriano de Competitividad.

%\textbf{Para esclarecer las causas sería necesario mirar el comportamiento de la inversión en cada rama de actividad y las innovaciones tecnológicas durante el periodo estudiado.} Adicionalmente, las variaciones en la productividad pueden también deberse a la reasignación del trabajo hacia actividades más (in)eficientes. Por ello, una agenda de productividad laboral debe también facilitar la reasignación de trabajadores hacia actividades de mayor productividad.

\textbf{A continuación se presenta la descomposición del crecimiento de cada una de las trece ramas comparables de actividad económica CIIU.} En particular, se presenta el PIB real de la actividad, el número de ocupados, el PIB por trabajador, las horas semanales por trabajador y el PIB por hora trabajada para los años 2010 y 2025. La lectura conjunta de estas variables permite distinguir si los aumentos en la producción responden principalmente a variaciones en el empleo, a cambios en las horas trabajadas o un aumento de la productividad por hora.

%Las trece ramas comparables parten de las agrupaciones reportadas por el DANE en cuentas nacionales. La diferencia frente a las doce agrupaciones usuales es que este informe separa la agrupación R+S+T en dos ramas: arte, entretenimiento, recreación y otros servicios, por un lado, y hogares como empleadores, por el otro. Esta separación es posible porque tanto las cuentas nacionales como la GEIH permiten construir numeradores y denominadores laborales comparables para ambas ramas. Como las series se expresan en pesos constantes, pueden presentarse pequeñas diferencias de aditividad entre niveles de agregación.

\subsection{Agropecuario}

\textbf{La actividad agropecuaria comprende cultivos, apoyo agropecuario y actividades mixtas; café; ganadería; silvicultura; y pesca y acuicultura.} En 2025 la actividad agropecuaria representó 6,3\% del PIB, con 63,9 billones de pesos de 2015, y concentró 13,6\% de los ocupados, con 3,1 millones de personas. El PIB por trabajador fue 20,4 millones de pesos de 2015 por ocupado, por debajo del agregado (44,5 millones). Las horas semanales por trabajador fueron 42,0, muy cerca del agregado (43,7). El PIB por hora trabajada fue 9,4 mil pesos de 2015 por hora, por debajo del agregado (19,6 mil).

\begin{table}[H]
\centering
\caption{Agropecuario: PIB, ocupados, horas y productividad laboral, 2010--2025}
\label{tab:sector_a_productividad}
\small
\begin{tabular}{lrrr}
\toprule
Indicador & 2010 & 2025 & Crec. anual \\
\midrule
PIB real  (Billones de pesos de 2015) & 39,9 & 63,9 & 3,18\% \\
Ocupados (Millones) & 3,0 & 3,1 & 0,38\% \\
PIB por trabajador (Millones de pesos de 2015) & 13,5 & 20,4 & 2,79\% \\
Horas semanales por trabajador & 45,0 & 42,0 & -0,46\% \\
PIB por hora trabajada (Miles de pesos de 2015) & 5,8 & 9,4 & 3,27\% \\
\bottomrule
\end{tabular}
\end{table}

\textbf{El crecimiento de la productividad del sector agropecuario fue superior al crecimiento de la productividad agregada.} Entre 2010 y 2025 el PIB de la actividad agropecuaria registró un crecimiento anualizado del 3,18\%, muy cerca del agregado (3,16\%). Los ocupados en la actividad registraron un crecimiento anual del 0,38\%, muy por debajo del agregado (2,09\%). El PIB por trabajador registró un crecimiento anual del 2,79\%, muy por encima del crecimiento del PIB por trabajador de toda la economía (1,05\%). Las horas semanales por trabajador registraron una caída anual del 0,46\%, muy cerca de la caída agregada (-0,45\%). El PIB por hora trabajada registró un crecimiento anual del 3,27\%, muy por encima del crecimiento del PIB por hora trabajada de toda la economía (1,51\%).

\begin{table}[H]
\centering
\caption{Productividad laboral de las actividades agropecuarias, 2010--2025}
\label{tab:sector_a_zoom61}
\small
\begin{tabular}{p{0.40\textwidth}rrrr}
\toprule
Actividad económica & PIB/trab. 2025 & Crec. & PIB/hora 2025 & Crec. \\
\midrule
Silvicultura & 70,2 & 5,15\% & 32,9 & 6,27\% \\
Pesca y acuicultura & 24,6 & 4,79\% & 10,9 & 5,75\% \\
Ganadería & 27,1 & 4,45\% & 12,0 & 5,10\% \\
Cultivos, apoyo agropecuario y mixtos & 20,5 & 2,24\% & 9,6 & 2,61\% \\
Café & 9,3 & 1,75\% & 4,1 & 1,98\% \\
\bottomrule
\end{tabular}
\caption*{\footnotesize Nota: PIB por trabajador en millones de pesos constantes de 2015; PIB por hora en miles de pesos constantes de 2015. %Cuando la GEIH no separa ocupados y horas al mismo nivel de las cuentas nacionales, se agrupan las subactividades del DANE hasta el nivel laboral comparable. 
Fuente: cálculos propios con DANE y GEIH.}
\end{table}

\textbf{El balance de la apertura del sector agropecuario es favorable: todas las subactividades superan el crecimiento agregado tanto en PIB por trabajador como en PIB por hora trabajada.} En esta apertura, el mayor crecimiento del PIB por trabajador se observa en Silvicultura (5,15\%), muy por encima del crecimiento agregado del PIB por trabajador (1,05\%); el menor se registra en Café (1,75\%), por encima del crecimiento agregado del PIB por trabajador (1,05\%). Por hora trabajada, el mejor desempeño corresponde a Silvicultura (6,27\%), muy por encima del crecimiento agregado del PIB por hora trabajada (1,51\%); el más rezagado corresponde a Café (1,98\%), por encima del crecimiento agregado del PIB por hora trabajada (1,51\%).

\subsection{Minas}

\textbf{La actividad de minas y canteras comprende petróleo, gas y apoyo conexo; otras minas y canteras; carbón; y minerales metalíferos.} En 2025 la actividad de minas y canteras representó 3,4\% del PIB, con 34,7 billones de pesos de 2015, y concentró 1,3\% de los ocupados, con 0,3 millones de personas. El PIB por trabajador fue 115,8 millones de pesos de 2015 por ocupado, por encima del agregado (44,5 millones). Las horas semanales por trabajador fueron 46,0, por encima del agregado (43,7). El PIB por hora trabajada fue 48,4 mil pesos de 2015 por hora, por encima del agregado (19,6 mil).

\begin{table}[H]
\centering
\caption{Minas: PIB, ocupados, horas y productividad laboral, 2010--2025}
\label{tab:sector_b_productividad}
\small
\begin{tabular}{lrrr}
\toprule
Indicador & 2010 & 2025 & Crec. anual \\
\midrule
PIB real  (Billones de pesos de 2015) & 38,4 & 34,7 & -0,67\% \\
Ocupados (Millones) & 0,2 & 0,3 & 3,04\% \\
PIB por trabajador (Millones de pesos de 2015) & 200,8 & 115,8 & -3,60\% \\
Horas semanales por trabajador & 49,7 & 46,0 & -0,52\% \\
PIB por hora trabajada (Miles de pesos de 2015) & 77,6 & 48,4 & -3,09\% \\
\bottomrule
\end{tabular}
\end{table}

\textbf{La productividad de minas y canteras cayó entre 2010 y 2025.} Entre 2010 y 2025 el PIB de la actividad de minas y canteras registró una caída anualizada del 0,67\%, muy por debajo del agregado (3,16\%). Los ocupados en la actividad registraron un crecimiento anual del 3,04\%, por encima del agregado (2,09\%). El PIB por trabajador registró una caída anual del 3,60\%, muy por debajo del crecimiento del PIB por trabajador de toda la economía (1,05\%). Las horas semanales por trabajador registraron una caída anual del 0,52\%, muy cerca de la caída agregada (-0,45\%). El PIB por hora trabajada registró una caída anual del 3,09\%, muy por debajo del crecimiento del PIB por hora trabajada de toda la economía (1,51\%).

\begin{table}[H]
\centering
\caption{Productividad laboral de las actividades mineras, 2010--2025}
\label{tab:sector_b_zoom61}
\small
\begin{tabular}{p{0.40\textwidth}rrrr}
\toprule
Actividad económica & PIB/trab. 2025 & Crec. & PIB/hora 2025 & Crec. \\
\midrule
Petróleo, gas y apoyo conexo & 576,1 & 0,07\% & 231,6 & 0,59\% \\
Otras minas y canteras & 45,6 & -0,89\% & 19,7 & -0,45\% \\
Carbón & 103,5 & -5,88\% & 41,8 & -5,33\% \\
Minerales metalíferos & 17,6 & -6,44\% & 7,5 & -5,98\% \\
\bottomrule
\end{tabular}
\caption*{\footnotesize Nota: PIB por trabajador en millones de pesos constantes de 2015; PIB por hora en miles de pesos constantes de 2015. %Cuando la GEIH no separa ocupados y horas al mismo nivel de las cuentas nacionales, se agrupan las subactividades del DANE hasta el nivel laboral comparable.
Fuente: cálculos propios con DANE y GEIH.}
\end{table}

\textbf{El balance de la apertura de la actividad de minas y canteras es negativo: ninguna subactividad supera el crecimiento agregado y tres de las cuatro subactividades caen tanto en PIB por trabajador como en PIB por hora trabajada.} En esta apertura, el mayor crecimiento del PIB por trabajador se observa en Petróleo, gas y apoyo conexo (0,07\%), por debajo del crecimiento agregado del PIB por trabajador (1,05\%); el menor se registra en Minerales metalíferos (-6,44\%), muy por debajo del crecimiento agregado del PIB por trabajador (1,05\%). Por hora trabajada, el mejor desempeño corresponde a Petróleo, gas y apoyo conexo (0,59\%), por debajo del crecimiento agregado del PIB por hora trabajada (1,51\%); el más rezagado corresponde a Minerales metalíferos (-5,98\%), muy por debajo del crecimiento agregado del PIB por hora trabajada (1,51\%).

\subsection{Manufactura}

\textbf{La actividad manufacturera agrupa alimentos, bebidas y tabaco; textiles, confecciones y cuero; madera, papel e impresión; refinación, químicos y minerales; metalurgia, maquinaria y equipo; y muebles y otras manufactureras.} En 2025 la actividad manufacturera representó 11,1\% del PIB, con 113,1 billones de pesos de 2015, y concentró 10,9\% de los ocupados, con 2,5 millones de personas. El PIB por trabajador fue 45,0 millones de pesos de 2015 por ocupado, muy cerca del agregado (44,5 millones). Las horas semanales por trabajador fueron 44,3, muy cerca del agregado (43,7). El PIB por hora trabajada fue 19,5 mil pesos de 2015 por hora, muy cerca del agregado (19,6 mil).

\begin{table}[H]
\centering
\caption{Manufactura: PIB, ocupados, horas y productividad laboral, 2010--2025}
\label{tab:sector_c_productividad}
\small
\begin{tabular}{lrrr}
\toprule
Indicador & 2010 & 2025 & Crec. anual \\
\midrule
PIB real  (Billones de pesos de 2015) & 88,0 & 113,1 & 1,69\% \\
Ocupados (Millones) & 2,2 & 2,5 & 0,91\% \\
PIB por trabajador (Millones de pesos de 2015) & 40,1 & 45,0 & 0,77\% \\
Horas semanales por trabajador & 46,7 & 44,3 & -0,35\% \\
PIB por hora trabajada (Miles de pesos de 2015) & 16,5 & 19,5 & 1,12\% \\
\bottomrule
\end{tabular}
\end{table}

\textbf{El crecimiento de la productividad de la manufactura fue inferior al crecimiento de la productividad agregada.} Entre 2010 y 2025 el PIB de la actividad manufacturera registró un crecimiento anualizado del 1,69\%, muy por debajo del agregado (3,16\%). Los ocupados en la actividad registraron un crecimiento anual del 0,91\%, muy por debajo del agregado (2,09\%). El PIB por trabajador registró un crecimiento anual del 0,77\%, por debajo del crecimiento del PIB por trabajador de toda la economía (1,05\%). Las horas semanales por trabajador registraron una caída anual del 0,35\%, muy cerca de la caída agregada (-0,45\%). El PIB por hora trabajada registró un crecimiento anual del 1,12\%, por debajo del crecimiento del PIB por hora trabajada de toda la economía (1,51\%).

\begin{table}[H]
\centering
\caption{Productividad laboral de las actividades manufactureras, 2010--2025}
\label{tab:sector_c_zoom25}
\small
\begin{tabular}{p{0.40\textwidth}rrrr}
\toprule
Actividad económica & PIB/trab. 2025 & Crec. & PIB/hora 2025 & Crec. \\
\midrule
Metalurgia, maquinaria y equipo & 48,8 & 1,69\% & 20,5 & 2,16\% \\
Alimentos, bebidas y tabaco & 50,8 & 1,60\% & 21,8 & 1,70\% \\
Madera, papel e impresión & 46,9 & 1,46\% & 21,8 & 2,21\% \\
Textiles, confecciones y cuero & 17,5 & 0,74\% & 7,8 & 1,19\% \\
Refinación, químicos y minerales & 106,6 & -0,75\% & 44,8 & -0,30\% \\
Muebles y otras manufactureras & 17,4 & -1,38\% & 7,6 & -1,13\% \\
\bottomrule
\end{tabular}
\caption*{\footnotesize Nota: PIB por trabajador en millones de pesos constantes de 2015; PIB por hora en miles de pesos constantes de 2015. Fuente: cálculos propios con DANE y GEIH.}
\end{table}

\textbf{El balance de la apertura de la manufactura es heterogéneo: tres de las seis subactividades superan el crecimiento agregado del PIB por trabajador y tres de las seis subactividades superan el crecimiento agregado del PIB por hora trabajada, mientras otras quedan rezagadas.} En esta apertura, el mayor crecimiento del PIB por trabajador se observa en Metalurgia, maquinaria y equipo (1,69\%), por encima del crecimiento agregado del PIB por trabajador (1,05\%); el menor se registra en Muebles y otras manufactureras (-1,38\%), muy por debajo del crecimiento agregado del PIB por trabajador (1,05\%). Por hora trabajada, el mejor desempeño corresponde a Madera, papel e impresión (2,21\%), por encima del crecimiento agregado del PIB por hora trabajada (1,51\%); el más rezagado corresponde a Muebles y otras manufactureras (-1,13\%), muy por debajo del crecimiento agregado del PIB por hora trabajada (1,51\%).

\subsection{Servicios públicos}

\textbf{La actividad de servicios públicos combina electricidad, gas y vapor con agua, saneamiento y desechos.} En 2025 la actividad de servicios públicos representó 3,0\% del PIB, con 30,3 billones de pesos de 2015, y concentró 1,3\% de los ocupados, con 0,3 millones de personas. El PIB por trabajador fue 99,5 millones de pesos de 2015 por ocupado, por encima del agregado (44,5 millones). Las horas semanales por trabajador fueron 44,0, muy cerca del agregado (43,7). El PIB por hora trabajada fue 43,5 mil pesos de 2015 por hora, por encima del agregado (19,6 mil).

\begin{table}[H]
\centering
\caption{Servicios públicos: PIB, ocupados, horas y productividad laboral, 2010--2025}
\label{tab:sector_d_e_productividad}
\small
\begin{tabular}{lrrr}
\toprule
Indicador & 2010 & 2025 & Crec. anual \\
\midrule
PIB real  (Billones de pesos de 2015) & 21,9 & 30,3 & 2,18\% \\
Ocupados (Millones) & 0,1 & 0,3 & 6,16\% \\
PIB por trabajador (Millones de pesos de 2015) & 176,5 & 99,5 & -3,75\% \\
Horas semanales por trabajador & 48,5 & 44,0 & -0,64\% \\
PIB por hora trabajada (Miles de pesos de 2015) & 70,0 & 43,5 & -3,13\% \\
\bottomrule
\end{tabular}
\end{table}

\textbf{La productividad de los servicios públicos cayó entre 2010 y 2025.} Entre 2010 y 2025 el PIB de la actividad de servicios públicos registró un crecimiento anualizado del 2,18\%, por debajo del agregado (3,16\%). Los ocupados en la actividad registraron un crecimiento anual del 6,16\%, muy por encima del agregado (2,09\%). El PIB por trabajador registró una caída anual del 3,75\%, muy por debajo del crecimiento del PIB por trabajador de toda la economía (1,05\%). Las horas semanales por trabajador registraron una caída anual del 0,64\%, muy cerca de la caída agregada (-0,45\%). El PIB por hora trabajada registró una caída anual del 3,13\%, muy por debajo del crecimiento del PIB por hora trabajada de toda la economía (1,51\%).

\begin{table}[H]
\centering
\caption{Productividad laboral de las actividades de servicios públicos, 2010--2025}
\label{tab:sector_d_e_zoom25}
\small
\begin{tabular}{p{0.40\textwidth}rrrr}
\toprule
Actividad económica & PIB/trab. 2025 & Crec. & PIB/hora 2025 & Crec. \\
\midrule
Electricidad, gas y vapor & 262,2 & -0,13\% & 109,6 & 0,28\% \\
Agua, saneamiento y desechos & 37,9 & -6,41\% & 16,9 & -5,74\% \\
\bottomrule
\end{tabular}
\caption*{\footnotesize Nota: PIB por trabajador en millones de pesos constantes de 2015; PIB por hora en miles de pesos constantes de 2015. Fuente: cálculos propios con DANE y GEIH.}
\end{table}

\textbf{El balance de la apertura de los servicios públicos es negativo: ninguna subactividad supera el crecimiento agregado, todas registran caídas en PIB por trabajador y una de las dos subactividades también cae por hora trabajada.} En esta apertura, el mayor crecimiento del PIB por trabajador se observa en Electricidad, gas y vapor (-0,13\%), muy por debajo del crecimiento agregado del PIB por trabajador (1,05\%); el menor se registra en Agua, saneamiento y desechos (-6,41\%), muy por debajo del crecimiento agregado del PIB por trabajador (1,05\%). Por hora trabajada, el mejor desempeño corresponde a Electricidad, gas y vapor (0,28\%), muy por debajo del crecimiento agregado del PIB por hora trabajada (1,51\%); el más rezagado corresponde a Agua, saneamiento y desechos (-5,74\%), muy por debajo del crecimiento agregado del PIB por hora trabajada (1,51\%).

\subsection{Construcción}

\textbf{La actividad de construcción comprende edificaciones, obras civiles y actividades especializadas de construcción.} En 2025 la actividad de construcción representó 4,0\% del PIB, con 41,3 billones de pesos de 2015, y concentró 6,9\% de los ocupados, con 1,6 millones de personas. El PIB por trabajador fue 26,1 millones de pesos de 2015 por ocupado, por debajo del agregado (44,5 millones). Las horas semanales por trabajador fueron 45,7, muy cerca del agregado (43,7). El PIB por hora trabajada fue 11,0 mil pesos de 2015 por hora, por debajo del agregado (19,6 mil).

\begin{table}[H]
\centering
\caption{Construcción: PIB, ocupados, horas y productividad laboral, 2010--2025}
\label{tab:sector_f_productividad}
\small
\begin{tabular}{lrrr}
\toprule
Indicador & 2010 & 2025 & Crec. anual \\
\midrule
PIB real  (Billones de pesos de 2015) & 40,0 & 41,3 & 0,21\% \\
Ocupados (Millones) & 1,0 & 1,6 & 3,39\% \\
PIB por trabajador (Millones de pesos de 2015) & 41,8 & 26,1 & -3,08\% \\
Horas semanales por trabajador & 48,9 & 45,7 & -0,46\% \\
PIB por hora trabajada (Miles de pesos de 2015) & 16,4 & 11,0 & -2,63\% \\
\bottomrule
\end{tabular}
\end{table}

\textbf{La productividad de la construcción cayó entre 2010 y 2025.} Entre 2010 y 2025 el PIB de la actividad de construcción registró un crecimiento anualizado del 0,21\%, muy por debajo del agregado (3,16\%). Los ocupados en la actividad registraron un crecimiento anual del 3,39\%, muy por encima del agregado (2,09\%). El PIB por trabajador registró una caída anual del 3,08\%, muy por debajo del crecimiento del PIB por trabajador de toda la economía (1,05\%). Las horas semanales por trabajador registraron una caída anual del 0,46\%, muy cerca de la caída agregada (-0,45\%). El PIB por hora trabajada registró una caída anual del 2,63\%, muy por debajo del crecimiento del PIB por hora trabajada de toda la economía (1,51\%).

\begin{table}[H]
\centering
\caption{Productividad laboral de las actividades de construcción, 2010--2025}
\label{tab:sector_f_zoom25}
\small
\begin{tabular}{p{0.40\textwidth}rrrr}
\toprule
Actividad económica & PIB/trab. 2025 & Crec. & PIB/hora 2025 & Crec. \\
\midrule
Obras civiles & 56,2 & -2,58\% & 23,1 & -2,01\% \\
Actividades especializadas de construcción & 26,7 & -3,27\% & 11,7 & -2,90\% \\
Edificaciones & 19,9 & -3,52\% & 8,3 & -3,08\% \\
\bottomrule
\end{tabular}
\caption*{\footnotesize Nota: PIB por trabajador en millones de pesos constantes de 2015; PIB por hora en miles de pesos constantes de 2015. Fuente: cálculos propios con DANE y GEIH.}
\end{table}

\textbf{El balance de la apertura de la construcción es claramente negativo: todas las subactividades registran caídas de productividad y ninguna alcanza el crecimiento agregado.} En esta apertura, el mayor crecimiento del PIB por trabajador se observa en Obras civiles (-2,58\%), muy por debajo del crecimiento agregado del PIB por trabajador (1,05\%); el menor se registra en Edificaciones (-3,52\%), muy por debajo del crecimiento agregado del PIB por trabajador (1,05\%). Por hora trabajada, el mejor desempeño corresponde a Obras civiles (-2,01\%), muy por debajo del crecimiento agregado del PIB por hora trabajada (1,51\%); el más rezagado corresponde a Edificaciones (-3,08\%), muy por debajo del crecimiento agregado del PIB por hora trabajada (1,51\%).

\subsection{Comercio, transporte y alojamiento}

\textbf{La agrupación de comercio, transporte y alojamiento combina comercio y reparación; transporte y almacenamiento; y alojamiento y servicios de comida.} En 2025 la agrupación de comercio, transporte y alojamiento representó 17,5\% del PIB, con 179,1 billones de pesos de 2015, y concentró 32,4\% de los ocupados, con 7,4 millones de personas. El PIB por trabajador fue 24,1 millones de pesos de 2015 por ocupado, por debajo del agregado (44,5 millones). Las horas semanales por trabajador fueron 45,8, muy cerca del agregado (43,7). El PIB por hora trabajada fue 10,1 mil pesos de 2015 por hora, por debajo del agregado (19,6 mil).

\begin{table}[H]
\centering
\caption{Comercio, transporte y alojamiento: PIB, ocupados, horas y productividad laboral, 2010--2025}
\label{tab:sector_g_h_i_productividad}
\small
\begin{tabular}{lrrr}
\toprule
Indicador & 2010 & 2025 & Crec. anual \\
\midrule
PIB real  (Billones de pesos de 2015) & 107,6 & 179,1 & 3,45\% \\
Ocupados (Millones) & 5,2 & 7,4 & 2,46\% \\
PIB por trabajador (Millones de pesos de 2015) & 20,9 & 24,1 & 0,97\% \\
Horas semanales por trabajador & 50,6 & 45,8 & -0,66\% \\
PIB por hora trabajada (Miles de pesos de 2015) & 7,9 & 10,1 & 1,64\% \\
\bottomrule
\end{tabular}
\end{table}

\textbf{La productividad de comercio, transporte y alojamiento evolucionó de forma parecida a la productividad agregada.} Entre 2010 y 2025 el PIB de la agrupación de comercio, transporte y alojamiento registró un crecimiento anualizado del 3,45\%, por encima del agregado (3,16\%). Los ocupados en la actividad registraron un crecimiento anual del 2,46\%, por encima del agregado (2,09\%). El PIB por trabajador registró un crecimiento anual del 0,97\%, muy cerca del crecimiento del PIB por trabajador de toda la economía (1,05\%). Las horas semanales por trabajador registraron una caída anual del 0,66\%, una caída más pronunciada que la del agregado (-0,45\%). El PIB por hora trabajada registró un crecimiento anual del 1,64\%, muy cerca del crecimiento del PIB por hora trabajada de toda la economía (1,51\%).

\begin{table}[H]
\centering
\caption{Productividad laboral de las actividades de comercio, transporte y alojamiento, 2010--2025}
\label{tab:sector_g_h_i_zoom25}
\small
\begin{tabular}{p{0.40\textwidth}rrrr}
\toprule
Actividad económica & PIB/trab. 2025 & Crec. & PIB/hora 2025 & Crec. \\
\midrule
Comercio y reparación & 22,8 & 2,28\% & 9,5 & 2,76\% \\
Transporte y almacenamiento & 27,8 & -0,18\% & 10,9 & 0,83\% \\
Alojamiento y comida & 23,4 & -1,08\% & 10,7 & -0,42\% \\
\bottomrule
\end{tabular}
\caption*{\footnotesize Nota: PIB por trabajador en millones de pesos constantes de 2015; PIB por hora en miles de pesos constantes de 2015. Fuente: cálculos propios con DANE y GEIH.}
\end{table}

\textbf{El balance de la apertura de la agrupación de comercio, transporte y alojamiento es heterogéneo: solo una de las tres subactividades supera en crecimiento de la productividad al resto de la economía.} En esta apertura, el mayor crecimiento del PIB por trabajador se observa en Comercio y reparación (2,28\%), muy por encima del crecimiento agregado del PIB por trabajador (1,05\%); el menor se registra en Alojamiento y comida (-1,08\%), muy por debajo del crecimiento agregado del PIB por trabajador (1,05\%). Por hora trabajada, el mejor desempeño corresponde a Comercio y reparación (2,76\%), muy por encima del crecimiento agregado del PIB por hora trabajada (1,51\%); el más rezagado corresponde a Alojamiento y comida (-0,42\%), muy por debajo del crecimiento agregado del PIB por hora trabajada (1,51\%).

\subsection{Información y comunicaciones}

\textbf{La actividad de información y comunicaciones reúne edición de contenidos audiovisuales, radio, televisión, telecomunicaciones, desarrollo de software, procesamiento de datos, y otros servicios de información.} En 2025 la actividad de información y comunicaciones representó 3,1\% del PIB, con 31,3 billones de pesos de 2015, y concentró 1,7\% de los ocupados, con 0,4 millones de personas. El PIB por trabajador fue 78,6 millones de pesos de 2015 por ocupado, por encima del agregado (44,5 millones). Las horas semanales por trabajador fueron 44,4, muy cerca del agregado (43,7). El PIB por hora trabajada fue 34,0 mil pesos de 2015 por hora, por encima del agregado (19,6 mil).

\begin{table}[H]
\centering
\caption{Información y comunicaciones: PIB, ocupados, horas y productividad laboral, 2010--2025}
\label{tab:sector_j_productividad}
\small
\begin{tabular}{lrrr}
\toprule
Indicador & 2010 & 2025 & Crec. anual \\
\midrule
PIB real  (Billones de pesos de 2015) & 18,3 & 31,3 & 3,66\% \\
Ocupados (Millones) & 0,4 & 0,4 & -0,07\% \\
PIB por trabajador (Millones de pesos de 2015) & 45,4 & 78,6 & 3,73\% \\
Horas semanales por trabajador & 49,6 & 44,4 & -0,73\% \\
PIB por hora trabajada (Miles de pesos de 2015) & 17,6 & 34,0 & 4,49\% \\
\bottomrule
\end{tabular}
\end{table}

\textbf{El crecimiento de la productividad de información y comunicaciones fue superior al crecimiento de la productividad agregada.} Entre 2010 y 2025 el PIB de la actividad de información y comunicaciones registró un crecimiento anualizado del 3,66\%, por encima del agregado (3,16\%). Los ocupados en la actividad registraron una caída anual del 0,07\%, muy por debajo del agregado (2,09\%). El PIB por trabajador registró un crecimiento anual del 3,73\%, muy por encima del crecimiento del PIB por trabajador de toda la economía (1,05\%). Las horas semanales por trabajador registraron una caída anual del 0,73\%, una caída más pronunciada que la del agregado (-0,45\%). El PIB por hora trabajada registró un crecimiento anual del 4,49\%, muy por encima del crecimiento del PIB por hora trabajada de toda la economía (1,51\%).

\subsection{Financieras}

\textbf{La actividad financiera reúne intermediación financiera y banca, seguros, fondos de pensiones y servicios auxiliares como administración de mercados, corretaje y otras actividades de apoyo financiero.} En 2025 la actividad financiera representó 5,2\% del PIB, con 53,1 billones de pesos de 2015, y concentró 1,9\% de los ocupados, con 0,4 millones de personas. El PIB por trabajador fue 124,5 millones de pesos de 2015 por ocupado, por encima del agregado (44,5 millones). Las horas semanales por trabajador fueron 44,2, muy cerca del agregado (43,7). El PIB por hora trabajada fue 54,2 mil pesos de 2015 por hora, por encima del agregado (19,6 mil).

\begin{table}[H]
\centering
\caption{Financieras: PIB, ocupados, horas y productividad laboral, 2010--2025}
\label{tab:sector_k_productividad}
\small
\begin{tabular}{lrrr}
\toprule
Indicador & 2010 & 2025 & Crec. anual \\
\midrule
PIB real  (Billones de pesos de 2015) & 22,3 & 53,1 & 5,95\% \\
Ocupados (Millones) & 0,2 & 0,4 & 4,49\% \\
PIB por trabajador (Millones de pesos de 2015) & 101,0 & 124,5 & 1,40\% \\
Horas semanales por trabajador & 45,2 & 44,2 & -0,15\% \\
PIB por hora trabajada (Miles de pesos de 2015) & 43,0 & 54,2 & 1,56\% \\
\bottomrule
\end{tabular}
\end{table}

\textbf{La evolución de la productividad de las actividades financieras fue parecida a la del resto de la economía.} Entre 2010 y 2025 el PIB de la actividad financiera registró un crecimiento anualizado del 5,95\%, muy por encima del agregado (3,16\%). Los ocupados en la actividad registraron un crecimiento anual del 4,49\%, muy por encima del agregado (2,09\%). El PIB por trabajador registró un crecimiento anual del 1,40\%, por encima del crecimiento del PIB por trabajador de toda la economía (1,05\%). Las horas semanales por trabajador registraron una caída anual del 0,15\%, una caída menos pronunciada que la del agregado (-0,45\%). El PIB por hora trabajada registró un crecimiento anual del 1,56\%, muy cerca del crecimiento del PIB por hora trabajada de toda la economía (1,51\%).

\subsection{Inmobiliarias}

\textbf{La actividad inmobiliaria reúne alquiler, compra, venta, administración e intermediación de bienes inmuebles, tanto propios o arrendados como de terceros.} En 2025 la actividad inmobiliaria representó 8,8\% del PIB, con 90,1 billones de pesos de 2015, y concentró 1,4\% de los ocupados, con 0,3 millones de personas. El PIB por trabajador fue 284,6 millones de pesos de 2015 por ocupado, por encima del agregado (44,5 millones). Las horas semanales por trabajador fueron 45,8, muy cerca del agregado (43,7). El PIB por hora trabajada fue 119,5 mil pesos de 2015 por hora, por encima del agregado (19,6 mil).

\begin{table}[H]
\centering
\caption{Inmobiliarias: PIB, ocupados, horas y productividad laboral, 2010--2025}
\label{tab:sector_l_productividad}
\small
\begin{tabular}{lrrr}
\toprule
Indicador & 2010 & 2025 & Crec. anual \\
\midrule
PIB real  (Billones de pesos de 2015) & 59,9 & 90,1 & 2,75\% \\
Ocupados (Millones) & 0,2 & 0,3 & 2,64\% \\
PIB por trabajador (Millones de pesos de 2015) & 280,0 & 284,6 & 0,11\% \\
Horas semanales por trabajador & 51,1 & 45,8 & -0,73\% \\
PIB por hora trabajada (Miles de pesos de 2015) & 105,3 & 119,5 & 0,85\% \\
\bottomrule
\end{tabular}
\end{table}

\textbf{El crecimiento de la productividad de las actividades inmobiliarias fue inferior al crecimiento de la productividad agregada.} Entre 2010 y 2025 el PIB de la actividad inmobiliaria registró un crecimiento anualizado del 2,75\%, por debajo del agregado (3,16\%). Los ocupados en la actividad registraron un crecimiento anual del 2,64\%, por encima del agregado (2,09\%). El PIB por trabajador registró un crecimiento anual del 0,11\%, por debajo del crecimiento del PIB por trabajador de toda la economía (1,05\%). Las horas semanales por trabajador registraron una caída anual del 0,73\%, una caída más pronunciada que la del agregado (-0,45\%). El PIB por hora trabajada registró un crecimiento anual del 0,85\%, por debajo del crecimiento del PIB por hora trabajada de toda la economía (1,51\%).

\subsection{Profesionales y administrativas}

\textbf{La agrupación de actividades profesionales y administrativas combina actividades profesionales, científicas y técnicas con servicios administrativos y de apoyo.} En 2025 la agrupación de actividades profesionales y administrativas representó 6,9\% del PIB, con 70,1 billones de pesos de 2015, y concentró 7,2\% de los ocupados, con 1,7 millones de personas. El PIB por trabajador fue 42,4 millones de pesos de 2015 por ocupado, muy cerca del agregado (44,5 millones). Las horas semanales por trabajador fueron 39,6, por debajo del agregado (43,7). El PIB por hora trabajada fue 20,6 mil pesos de 2015 por hora, por encima del agregado (19,6 mil).

\begin{table}[H]
\centering
\caption{Profesionales y administrativas: PIB, ocupados, horas y productividad laboral, 2010--2025}
\label{tab:sector_m_n_productividad}
\small
\begin{tabular}{lrrr}
\toprule
Indicador & 2010 & 2025 & Crec. anual \\
\midrule
PIB real  (Billones de pesos de 2015) & 45,4 & 70,1 & 2,94\% \\
Ocupados (Millones) & 0,8 & 1,7 & 4,65\% \\
PIB por trabajador (Millones de pesos de 2015) & 54,3 & 42,4 & -1,63\% \\
Horas semanales por trabajador & 42,1 & 39,6 & -0,41\% \\
PIB por hora trabajada (Miles de pesos de 2015) & 24,8 & 20,6 & -1,22\% \\
\bottomrule
\end{tabular}
\end{table}

\textbf{La productividad de las actividades profesionales y administrativas cayó entre 2010 y 2025.} Entre 2010 y 2025 el PIB de la agrupación de actividades profesionales y administrativas registró un crecimiento anualizado del 2,94\%, por debajo del agregado (3,16\%). Los ocupados en la actividad registraron un crecimiento anual del 4,65\%, muy por encima del agregado (2,09\%). El PIB por trabajador registró una caída anual del 1,63\%, muy por debajo del crecimiento del PIB por trabajador de toda la economía (1,05\%). Las horas semanales por trabajador registraron una caída anual del 0,41\%, muy cerca de la caída agregada (-0,45\%). El PIB por hora trabajada registró una caída anual del 1,22\%, muy por debajo del crecimiento del PIB por hora trabajada de toda la economía (1,51\%).

\begin{table}[H]
\centering
\caption{Productividad laboral de las actividades profesionales y administrativas, 2010--2025}
\label{tab:sector_m_n_zoom61}
\small
\begin{tabular}{p{0.40\textwidth}rrrr}
\toprule
Actividad económica & PIB/trab. 2025 & Crec. & PIB/hora 2025 & Crec. \\
\midrule
Servicios administrativos y de apoyo & 36,0 & -1,32\% & 18,0 & -0,87\% \\
Profesionales, científicas y técnicas & 54,7 & -1,91\% & 25,2 & -1,60\% \\
\bottomrule
\end{tabular}
\caption*{\footnotesize Nota: PIB por trabajador en millones de pesos constantes de 2015; PIB por hora en miles de pesos constantes de 2015. %Cuando la GEIH no separa ocupados y horas al mismo nivel de las cuentas nacionales, se agrupan las subactividades del DANE hasta el nivel laboral comparable.
Fuente: cálculos propios con DANE y GEIH.}
\end{table}

\textbf{El balance de la apertura de las actividades profesionales y administrativas es claramente negativo: todas las subactividades registran caídas de productividad y ninguna alcanza el crecimiento agregado.} En esta apertura, el mayor crecimiento del PIB por trabajador se observa en Servicios administrativos y de apoyo (-1,32\%), muy por debajo del crecimiento agregado del PIB por trabajador (1,05\%); el menor se registra en Profesionales, científicas y técnicas (-1,91\%), muy por debajo del crecimiento agregado del PIB por trabajador (1,05\%). Por hora trabajada, el mejor desempeño corresponde a Servicios administrativos y de apoyo (-0,87\%), muy por debajo del crecimiento agregado del PIB por hora trabajada (1,51\%); el más rezagado corresponde a Profesionales, científicas y técnicas (-1,60\%), muy por debajo del crecimiento agregado del PIB por hora trabajada (1,51\%).

\subsection{Adm. pública, educación y salud}

\textbf{La agrupación de administración pública, educación y salud combina administración pública y defensa, educación, salud humana y servicios sociales.} En 2025 la agrupación de administración pública, educación y salud representó 16,2\% del PIB, con 165,3 billones de pesos de 2015, y concentró 12,3\% de los ocupados, con 2,8 millones de personas. El PIB por trabajador fue 58,4 millones de pesos de 2015 por ocupado, por encima del agregado (44,5 millones). Las horas semanales por trabajador fueron 43,0, muy cerca del agregado (43,7). El PIB por hora trabajada fue 26,1 mil pesos de 2015 por hora, por encima del agregado (19,6 mil).

\begin{table}[H]
\centering
\caption{Adm. pública, educación y salud: PIB, ocupados, horas y productividad laboral, 2010--2025}
\label{tab:sector_o_p_q_productividad}
\small
\begin{tabular}{lrrr}
\toprule
Indicador & 2010 & 2025 & Crec. anual \\
\midrule
PIB real  (Billones de pesos de 2015) & 85,4 & 165,3 & 4,51\% \\
Ocupados (Millones) & 1,9 & 2,8 & 2,54\% \\
PIB por trabajador (Millones de pesos de 2015) & 43,9 & 58,4 & 1,92\% \\
Horas semanales por trabajador & 43,3 & 43,0 & -0,03\% \\
PIB por hora trabajada (Miles de pesos de 2015) & 19,5 & 26,1 & 1,95\% \\
\bottomrule
\end{tabular}
\end{table}

\textbf{El crecimiento de la productividad de administración pública, educación y salud fue superior al crecimiento de la productividad agregada.} Entre 2010 y 2025 el PIB de la agrupación de administración pública, educación y salud registró un crecimiento anualizado del 4,51\%, muy por encima del agregado (3,16\%). Los ocupados en la actividad registraron un crecimiento anual del 2,54\%, por encima del agregado (2,09\%). El PIB por trabajador registró un crecimiento anual del 1,92\%, por encima del crecimiento del PIB por trabajador de toda la economía (1,05\%). Las horas semanales por trabajador registraron una caída anual del 0,03\%, una caída menos pronunciada que la del agregado (-0,45\%). El PIB por hora trabajada registró un crecimiento anual del 1,95\%, por encima del crecimiento del PIB por hora trabajada de toda la economía (1,51\%).

\begin{table}[H]
\centering
\caption{Productividad laboral de las actividades de administración pública, educación y salud, 2010--2025}
\label{tab:sector_o_p_q_zoom25}
\small
\begin{tabular}{p{0.40\textwidth}rrrr}
\toprule
Actividad económica & PIB/trab. 2025 & Crec. & PIB/hora 2025 & Crec. \\
\midrule
Salud y servicios sociales & 49,1 & 2,54\% & 21,3 & 2,47\% \\
Educación & 52,3 & 1,54\% & 25,8 & 1,26\% \\
Administración pública & 77,2 & 1,40\% & 32,2 & 2,06\% \\
\bottomrule
\end{tabular}
\caption*{\footnotesize Nota: PIB por trabajador en millones de pesos constantes de 2015; PIB por hora en miles de pesos constantes de 2015. Fuente: cálculos propios con DANE y GEIH.}
\end{table}

\textbf{El balance de la apertura de administración pública, educación y salud es favorable en PIB por trabajador, pero menos uniforme por hora trabajada: todas las subactividades superan el crecimiento agregado por trabajador y dos de las tres subactividades lo hacen por hora.} En esta apertura, el mayor crecimiento del PIB por trabajador se observa en Salud y servicios sociales (2,54\%), muy por encima del crecimiento agregado del PIB por trabajador (1,05\%); el menor se registra en Administración pública (1,40\%), por encima del crecimiento agregado del PIB por trabajador (1,05\%). Por hora trabajada, el mejor desempeño corresponde a Salud y servicios sociales (2,47\%), por encima del crecimiento agregado del PIB por hora trabajada (1,51\%); el más rezagado corresponde a Educación (1,26\%), por debajo del crecimiento agregado del PIB por hora trabajada (1,51\%).

\subsection{Arte, entretenimiento, recreación y otros servicios}

\textbf{La agrupación de arte, entretenimiento, recreación y otros servicios reúne actividades artísticas, de entretenimiento y recreación, juegos de azar, actividades deportivas, servicios personales, asociaciones y reparación de computadores y efectos personales.} En 2025 la agrupación de arte, entretenimiento, recreación y otros servicios representó 4,0\% del PIB, con 40,6 billones de pesos de 2015, y concentró 5,4\% de los ocupados, con 1,2 millones de personas. El PIB por trabajador fue 32,8 millones de pesos de 2015 por ocupado, por debajo del agregado (44,5 millones). Las horas semanales por trabajador fueron 39,6, por debajo del agregado (43,7). El PIB por hora trabajada fue 16,0 mil pesos de 2015 por hora, por debajo del agregado (19,6 mil).

\begin{table}[H]
\centering
\caption{Arte, entretenimiento, recreación y otros servicios: PIB, ocupados, horas y productividad laboral, 2010--2025}
\label{tab:sector_r_s_productividad}
\small
\begin{tabular}{lrrr}
\toprule
Indicador & 2010 & 2025 & Crec. anual \\
\midrule
PIB real  (Billones de pesos de 2015) & 11,0 & 40,6 & 9,09\% \\
Ocupados (Millones) & 1,0 & 1,2 & 1,65\% \\
PIB por trabajador (Millones de pesos de 2015) & 11,4 & 32,8 & 7,32\% \\
Horas semanales por trabajador & 39,3 & 39,6 & 0,04\% \\
PIB por hora trabajada (Miles de pesos de 2015) & 5,6 & 16,0 & 7,28\% \\
\bottomrule
\end{tabular}
\end{table}

\textbf{El crecimiento de la productividad de arte, entretenimiento, recreación y otros servicios fue superior al crecimiento de la productividad agregada.} Entre 2010 y 2025 el PIB de la agrupación de arte, entretenimiento, recreación y otros servicios registró un crecimiento anualizado del 9,09\%, muy por encima del agregado (3,16\%). Los ocupados en la actividad registraron un crecimiento anual del 1,65\%, por debajo del agregado (2,09\%). El PIB por trabajador registró un crecimiento anual del 7,32\%, muy por encima del crecimiento del PIB por trabajador de toda la economía (1,05\%). Las horas semanales por trabajador registraron un crecimiento anual del 0,04\%, por encima del agregado (-0,45\%). El PIB por hora trabajada registró un crecimiento anual del 7,28\%, muy por encima del crecimiento del PIB por hora trabajada de toda la economía (1,51\%).

\subsection{Hogares como empleadores}

\textbf{La actividad de hogares como empleadores reúne el trabajo doméstico remunerado en hogares y la producción no diferenciada de los hogares para uso propio.} En 2025 la actividad de hogares como empleadores representó 0,7\% del PIB, con 6,7 billones de pesos de 2015, y concentró 3,1\% de los ocupados, con 0,7 millones de personas. El PIB por trabajador fue 9,4 millones de pesos de 2015 por ocupado, por debajo del agregado (44,5 millones). Las horas semanales por trabajador fueron 40,5, por debajo del agregado (43,7). El PIB por hora trabajada fue 4,5 mil pesos de 2015 por hora, por debajo del agregado (19,6 mil).

\begin{table}[H]
\centering
\caption{Hogares como empleadores: PIB, ocupados, horas y productividad laboral, 2010--2025}
\label{tab:sector_t_productividad}
\small
\begin{tabular}{lrrr}
\toprule
Indicador & 2010 & 2025 & Crec. anual \\
\midrule
PIB real  (Billones de pesos de 2015) & 4,3 & 6,7 & 3,06\% \\
Ocupados (Millones) & 0,7 & 0,7 & 0,42\% \\
PIB por trabajador (Millones de pesos de 2015) & 6,4 & 9,4 & 2,63\% \\
Horas semanales por trabajador & 45,9 & 40,5 & -0,83\% \\
PIB por hora trabajada (Miles de pesos de 2015) & 2,7 & 4,5 & 3,49\% \\
\bottomrule
\end{tabular}
\end{table}

\textbf{El crecimiento de la productividad de hogares como empleadores fue superior al crecimiento de la productividad agregada.} Entre 2010 y 2025 el PIB de la actividad de hogares como empleadores registró un crecimiento anualizado del 3,06\%, muy cerca del agregado (3,16\%). Los ocupados en la actividad registraron un crecimiento anual del 0,42\%, muy por debajo del agregado (2,09\%). El PIB por trabajador registró un crecimiento anual del 2,63\%, muy por encima del crecimiento del PIB por trabajador de toda la economía (1,05\%). Las horas semanales por trabajador registraron una caída anual del 0,83\%, una caída más pronunciada que la del agregado (-0,45\%). El PIB por hora trabajada registró un crecimiento anual del 3,49\%, muy por encima del crecimiento del PIB por hora trabajada de toda la economía (1,51\%).

%{\footnotesize Nota general: para facilitar la lectura, el informe usa PIB por actividad económica; en sentido estricto, el numerador corresponde al valor agregado bruto de cada actividad reportado por el DANE. Valores en pesos constantes de 2015. Ocupados expandidos con el factor \texttt{fex}. Las horas semanales promedio se calculan como horas anuales totales divididas por ocupados y por 52 semanas. Se excluye 2020 por no contar con GEIH anual comparable en la base del proyecto. Fuente: cálculos propios con DANE y GEIH.}


\section{Crecimiento de la productividad y crecimiento de la ocupación}

\textbf{Después de revisar el crecimiento de la productividad laboral por actividades económicas, conviene contrastarlo con el crecimiento de la ocupación por actividad económica.} El apéndice de este informe presenta los indicadores de productividad y ocupación de 50 subactividades que se relacionan con las trece ramas comparables de actividad económica. Esta sección usa los resultados presentados en el apéndice para hacer dos preguntas. La primera es si las actividades donde más crecieron los ocupados y las horas totales fueron también aquellas donde más aumentó la productividad laboral. La segunda es si las actividades que partían de menores niveles de productividad en 2010 crecieron más rápido durante los años siguientes.

Para responder la primera pregunta, la Figura \ref{fig:pib_geih_productividad_sector_correlaciones} muestra los diagramas de dispersión correspondientes al crecimiento del PIB por hora / PIB por trabajador versus el crecimiento de los ocupados/horas trabajadas a través de 50 subactividades. En cada panel, \(r\) corresponde a la correlación entre las dos variables.

\begin{figure}[H]
  \centering
  \includegraphics[width=\textwidth]{Paper/figures/fig_pib_geih_productividad_sector_correlaciones.png}
  \caption{Crecimiento de la ocupación y versus crecimiento de la productividad}
  \label{fig:pib_geih_productividad_sector_correlaciones}
\end{figure}

\textbf{Las actividades donde más crecen los ocupados y las horas trabajadas no son las mismas donde más aumenta la productividad.} La correlación entre el crecimiento  de ocupados /horas  y el crecimiento del PIB por trabajador/hora oscila entre es igual cercana a -0,85. Esto no implica que la absorción de trabajadores se ha concentrado parcialmente en actividades donde el PIB por trabajador crece menos. 

\begin{figure}[H]
  \centering
  \includegraphics[width=\textwidth]{Paper/figures/fig_pib_geih_productividad_nivel_inicial_crecimiento.png}
  \caption{Nivel inicial de productividad y crecimiento posterior por actividad económica comparable}
  \label{fig:pib_geih_productividad_nivel_inicial_crecimiento}
\end{figure}

La Figura \ref{fig:pib_geih_productividad_nivel_inicial_crecimiento} agrega una lectura complementaria. La relación negativa entre el nivel inicial de productividad y su crecimiento posterior sugiere un patrón de convergencia parcial entre actividades económicas: varias de las actividades que partían de niveles más bajos en 2010 crecieron más rápido entre 2010 y 2025. En otras palabras, las actividades con menor productividad inicial crecieron relativamente más rápido.

\textbf{Los anteriores resultados muestran que aunque el empleo fue absorbido por actividades que eventualmente crecieron menos rápido, estas actividades contaban con una productividad inicial relativamente alta.} Por lo tanto los resultados deben tomarse con cautela, ya que no representan una mala asignación del trabajo. Sin embargo, no deja de ser cierto que el crecimiento del PIB hubiese sido mayor si las actividades que absorbieron el empleo hubiesen tenido un crecimiento de la productividad similar al que tuvieron actividades con una productividad inicial más baja.

\section{Conclusiones}

\textbf{Colombia presenta una brecha entre su productividad medida a través del producto por trabajador y su productividad medida a través del producto por hora.} Entre 2010 y 2025, el PIB por trabajador creció cerca de 1\% anual, mientras que el PIB por hora trabajada creció al 1,5\% anual. Esta brecha se explica por la caída anual de casi 0,5\% en las horas trabajadas por persona ocupada. Concentrarse exclusivamente en el producto por trabajador pasaría por alto que una economía también puede mejorar el bienestar si logra producir más por hora y, al mismo tiempo, liberar tiempo de trabajo, siempre que esa reducción de horas no refleje desempleo involuntario.

\textbf{La mirada por actividad económica muestra una gran diversidad en el crecimiento de la productividad.} Algunas actividades  --como arte, entretenimiento, recreación y otros servicios e información y comunicaciones-- registraron crecimientos de la productividad muy por encima del resto de la economía. Otras en cambio como la minería y la construcción presentaron caídas en el PIB por trabajador y en el PIB por hora. Además, las actividades que más aumentaron sus ocupados no fueron necesariamente las que más elevaron su productividad.

\textbf{La agenda de productividad para Colombia debe identificar las razones de este comportamiento tan heterogéneo.} Si bien las diferentes velocidades a las que crece la productividad de las actividades económicas pueden explicarse en parte por factores tecnológicos exógenos, también es importante averiguar cómo las decisiones de inversión han afectado la productividad laboral. En la medida en que la inversión se contrae, la productividad laboral puede disminuir como consecuencia de una menor disponibilidad de capital por trabajador. Sin embargo, el crecimiento de la ocupación en actividades con bajo crecimiento de la productividad también puede señalar que la economía colombiana genera incentivos a orientar los factores de producción hacia actividades con menor dinamismo productivo. Aunque el alcance del informe no permite corroborar esta hipótesis, la agenda del país debe identificar si este es el caso, con el fin de alinear los incentivos de manera más acorde con la asignación eficiente del trabajo. De esta manera, Colombia podría obtener un crecimiento más acelerado de la productividad y, con ello, un mayor ingreso por habitante, una mayor disponibilidad de tiempo libre, o ambas cosas.


\appendix
\section*{Apéndice}
\addcontentsline{toc}{section}{Apéndice}

%\subsection*{Descomposición del crecimiento por actividad económica}

\begin{figure}[H]
  \centering
  \makebox[\textwidth][c]{\includegraphics[width=1.03\textwidth]{Paper/figures/fig_pib_geih_descomposicion_sector.png}}
  \caption{Descomposición del crecimiento anual del PIB por actividad económica, 2010--2025}
  \label{fig:pib_geih_descomposicion_sector}
  \caption*{\footnotesize %Nota: las barras intermedias muestran contribuciones aproximadas en puntos porcentuales al crecimiento anualizado del PIB real de cada actividad económica. La barra final muestra el crecimiento anualizado del PIB de la actividad. A nivel de actividad económica, el numerador corresponde al valor agregado bruto. Por redondeo, las contribuciones pueden no sumar exactamente al crecimiento reportado. 
  Fuente: cálculos propios con base en Cuentas Nacionales y GEIH del DANE.}
\end{figure}

%\subsection*{Productividad por subactividades}

%La Figura \ref{fig:pib_geih_productividad_sector_indices} muestra la evolución del PIB por trabajador y del PIB por hora trabajada para las trece grandes ramas comparables de actividad económica. Ambas series se presentan como índices con base 2010 = 100, lo que permite comparar la trayectoria de las dos medidas de productividad dentro de cada actividad.

\begin{figure}[H]
  \centering
  \makebox[\textwidth][c]{\includegraphics[width=1.03\textwidth]{Paper/figures/fig_pib_geih_productividad_sector_indices.png}}
  \caption{Productividad laboral por grandes ramas de actividad económica, índice 2010 = 100}
  \label{fig:pib_geih_productividad_sector_indices}
  \caption*{\footnotesize Nota: el PIB por trabajador y el PIB por hora trabajada se indexan a 2010 = 100. Fuente: cálculos propios con DANE y GEIH.}
\end{figure}

%El Cuadro \ref{tab:pib_geih_productividad_61} presenta una desagregación de las actividades económicas mayor a la del Cuadro \ref{tab:pib_geih_productividad_sector}. Cuando la GEIH no permite separar ocupados y horas con el mismo detalle, las subactividades del DANE aparecen agrupadas en la observación laboral comparable.


\begingroup
\footnotesize
\begin{longtable}{p{0.34\textwidth}rrrrrr}
\caption{Productividad laboral por subactividades, 2010--2025}\label{tab:pib_geih_productividad_61}\\
\toprule
& \multicolumn{3}{c}{PIB por trabajador} & \multicolumn{3}{c}{PIB por hora} \\
& \multicolumn{3}{c}{\footnotesize Millones de pesos de 2015} & \multicolumn{3}{c}{\footnotesize Miles de pesos de 2015} \\
\cmidrule(lr){2-4}\cmidrule(lr){5-7}
Actividad económica & 2010 & 2025 & Crec. & 2010 & 2025 & Crec. \\
\midrule
\endfirsthead
\toprule
& \multicolumn{3}{c}{PIB por trabajador} & \multicolumn{3}{c}{PIB por hora} \\
& \multicolumn{3}{c}{\footnotesize Millones de pesos de 2015} & \multicolumn{3}{c}{\footnotesize Miles de pesos de 2015} \\
\cmidrule(lr){2-4}\cmidrule(lr){5-7}
Actividad económica & 2010 & 2025 & Crec. & 2010 & 2025 & Crec. \\
\midrule
\endhead
Arte, recreación y otros servicios & 11,4 & 32,8 & 7,32\% & 5,6 & 16,0 & 7,28\% \\
Silvicultura & 33,1 & 70,2 & 5,15\% & 13,2 & 32,9 & 6,27\% \\
Pesca y acuicultura & 12,2 & 24,6 & 4,79\% & 4,7 & 10,9 & 5,75\% \\
Ganadería & 14,1 & 27,1 & 4,45\% & 5,7 & 12,0 & 5,10\% \\
Información y comunicaciones & 45,4 & 78,6 & 3,73\% & 17,6 & 34,0 & 4,49\% \\
Transporte aéreo & 183,2 & 286,1 & 3,02\% & 72,9 & 126,4 & 3,74\% \\
Impresión y reproducción & 29,4 & 44,8 & 2,85\% & 12,2 & 20,1 & 3,40\% \\
Hogares como empleadores & 6,4 & 9,4 & 2,63\% & 2,7 & 4,5 & 3,49\% \\
Salud y servicios sociales & 33,7 & 49,1 & 2,54\% & 14,7 & 21,3 & 2,47\% \\
Equipo eléctrico y electrónico & 39,0 & 56,0 & 2,45\% & 15,5 & 23,6 & 2,83\% \\
Comercio; Reparación de vehículos & 16,3 & 22,8 & 2,28\% & 6,3 & 9,5 & 2,76\% \\
Cultivos, apoyo agropecuario y mixtos & 14,7 & 20,5 & 2,24\% & 6,5 & 9,6 & 2,61\% \\
Almacenamiento y apoyo al transporte & 27,6 & 37,2 & 2,02\% & 10,7 & 15,8 & 2,59\% \\
Café & 7,2 & 9,3 & 1,75\% & 3,1 & 4,1 & 1,98\% \\
Textiles y confecciones & 16,5 & 21,3 & 1,70\% & 7,1 & 9,5 & 1,91\% \\
Vehículos y equipo de transporte & 48,5 & 62,4 & 1,70\% & 18,6 & 26,2 & 2,31\% \\
Alimentos, bebidas y tabaco & 40,1 & 50,9 & 1,61\% & 17,0 & 21,9 & 1,71\% \\
Educación de mercado; Educación de no mercado & 41,6 & 52,4 & 1,56\% & 21,4 & 25,8 & 1,27\% \\
Otras manufactureras & 32,8 & 41,2 & 1,53\% & 15,9 & 20,6 & 1,71\% \\
Financieras y seguros & 101,0 & 124,5 & 1,40\% & 43,0 & 54,2 & 1,56\% \\
Administración pública y defensa & 62,7 & 77,2 & 1,40\% & 23,7 & 32,2 & 2,06\% \\
Papel y cartón & 77,3 & 89,1 & 0,96\% & 30,6 & 37,6 & 1,37\% \\
Inmobiliarias & 280,0 & 284,6 & 0,11\% & 105,3 & 119,5 & 0,85\% \\
Petróleo, gas y apoyo conexo & 570,5 & 576,1 & 0,07\% & 211,9 & 231,6 & 0,59\% \\
Gas, vapor y aire acondicionado & 155,7 & 154,9 & -0,03\% & 61,4 & 64,5 & 0,33\% \\
Muebles y colchones & 13,8 & 13,6 & -0,09\% & 5,3 & 5,6 & 0,35\% \\
Energía eléctrica & 303,3 & 298,7 & -0,10\% & 118,9 & 125,0 & 0,33\% \\
Metalurgia y productos metálicos & 31,5 & 30,5 & -0,22\% & 12,3 & 12,8 & 0,25\% \\
Madera & 19,5 & 18,5 & -0,38\% & 8,5 & 9,6 & 0,81\% \\
Minerales no metálicos & 124,7 & 113,9 & -0,60\% & 47,8 & 47,8 & -0,01\% \\
Transporte terrestre y tuberías & 27,5 & 24,5 & -0,78\% & 8,9 & 9,3 & 0,33\% \\
Químicos y farmacéuticos & 97,7 & 86,3 & -0,82\% & 39,5 & 36,7 & -0,49\% \\
Otras minas y canteras & 52,2 & 45,6 & -0,89\% & 21,1 & 19,7 & -0,45\% \\
Coquización y refinación & 2.441,9 & 2.103,3 & -0,99\% & 933,9 & 882,8 & -0,37\% \\
Alojamiento y comida & 27,6 & 23,4 & -1,08\% & 11,4 & 10,7 & -0,42\% \\
Maquinaria e instalación & 44,8 & 38,1 & -1,09\% & 17,8 & 16,5 & -0,50\% \\
Caucho y plástico & 34,4 & 28,3 & -1,28\% & 13,3 & 11,7 & -0,83\% \\
Servicios administrativos y de apoyo & 43,9 & 36,0 & -1,32\% & 20,5 & 18,0 & -0,87\% \\
Profesionales, científicas y técnicas & 73,0 & 54,7 & -1,91\% & 32,1 & 25,2 & -1,60\% \\
Transporte acuático & 16,1 & 11,0 & -2,52\% & 6,0 & 4,6 & -1,76\% \\
Obras civiles & 83,2 & 56,2 & -2,58\% & 31,3 & 23,1 & -2,01\% \\
Aguas residuales, desechos y saneamiento & 118,8 & 73,3 & -3,17\% & 46,7 & 31,2 & -2,64\% \\
Agua & 61,1 & 37,6 & -3,19\% & 24,7 & 15,6 & -3,02\% \\
Actividades especializadas de construcción & 44,0 & 26,7 & -3,27\% & 18,3 & 11,7 & -2,90\% \\
Edificaciones & 34,1 & 19,9 & -3,52\% & 13,3 & 8,3 & -3,08\% \\
Correo y mensajería & 11,2 & 6,4 & -3,61\% & 4,5 & 2,7 & -3,41\% \\
Cuero, calzado y artículos de viaje & 12,5 & 6,5 & -4,31\% & 4,7 & 2,9 & -3,14\% \\
Carbón & 256,7 & 103,5 & -5,88\% & 95,1 & 41,8 & -5,33\% \\
Minerales metalíferos & 47,7 & 17,6 & -6,44\% & 18,8 & 7,5 & -5,98\% \\
Reciclaje & 462,1 & 14,9 & -20,45\% & 197,5 & 7,2 & -19,83\% \\
\bottomrule
\end{longtable}
\endgroup
{\footnotesize Nota:  A nivel de actividad económica, el numerador de los indicadores de productividad corresponde estrictamente al valor agregado bruto de cada actividad. El renglón Alimentos, bebidas y tabaco agrega manufacturas de carnes y pescado, aceites, lácteos, molinería y panadería, café, azúcar y panela, cacao y confitería, frutas y hortalizas, otros alimentos, bebidas y tabaco. Fuente: cálculos propios con DANE y GEIH.}


%\subsection*{PIB y ocupados por subactividad}
\begingroup
\footnotesize
\begin{longtable}{p{0.34\textwidth}rrrrrr}
\caption{PIB y ocupados por subactividades, 2010--2025}\label{tab:pib_ocupados_61_comparable}\\
\toprule
& \multicolumn{3}{c}{PIB real} & \multicolumn{3}{c}{Ocupados} \\
& \multicolumn{3}{c}{\footnotesize Billones de pesos de 2015} & \multicolumn{3}{c}{\footnotesize Millones de personas} \\
\cmidrule(lr){2-4}\cmidrule(lr){5-7}
Actividad económica & 2010 & 2025 & Crec. anual & 2010 & 2025 & Crec. anual \\
\midrule
\endfirsthead
\toprule
& \multicolumn{3}{c}{PIB real} & \multicolumn{3}{c}{Ocupados} \\
& \multicolumn{3}{c}{\footnotesize Billones de pesos de 2015} & \multicolumn{3}{c}{\footnotesize Millones de personas} \\
\cmidrule(lr){2-4}\cmidrule(lr){5-7}
Actividad económica & 2010 & 2025 & Crec. anual & 2010 & 2025 & Crec. anual \\
\midrule
\endhead
Cultivos, apoyo agropecuario y mixtos & 22,9 & 37,2 & 3,31\% & 1,55 & 1,82 & 1,05\% \\
Café & 3,4 & 5,0 & 2,58\% & 0,47 & 0,54 & 0,81\% \\
Ganadería & 10,9 & 17,3 & 3,13\% & 0,77 & 0,64 & -1,26\% \\
Silvicultura & 1,4 & 2,1 & 2,75\% & 0,04 & 0,03 & -2,28\% \\
Pesca y acuicultura & 1,4 & 2,6 & 4,44\% & 0,11 & 0,11 & -0,34\% \\
Carbón & 7,7 & 5,7 & -1,98\% & 0,03 & 0,06 & 4,14\% \\
Petróleo, gas y apoyo conexo & 25,4 & 24,6 & -0,22\% & 0,04 & 0,04 & -0,29\% \\
Minerales metalíferos & 4,2 & 3,0 & -2,32\% & 0,09 & 0,17 & 4,40\% \\
Otras minas y canteras & 1,5 & 1,5 & 0,01\% & 0,03 & 0,03 & 0,91\% \\
Alimentos, bebidas y tabaco & 24,0 & 33,5 & 2,27\% & 0,60 & 0,66 & 0,65\% \\
Textiles y confecciones & 8,0 & 10,6 & 1,93\% & 0,48 & 0,50 & 0,23\% \\
Cuero, calzado y artículos de viaje & 1,7 & 1,2 & -2,28\% & 0,13 & 0,18 & 2,11\% \\
Madera & 1,1 & 1,0 & -0,26\% & 0,05 & 0,05 & 0,12\% \\
Papel y cartón & 2,8 & 3,6 & 1,77\% & 0,04 & 0,04 & 0,81\% \\
Impresión y reproducción & 1,6 & 1,9 & 1,17\% & 0,05 & 0,04 & -1,63\% \\
Coquización y refinación & 8,6 & 11,4 & 1,92\% & 0,00 & 0,01 & 2,94\% \\
Químicos y farmacéuticos & 10,2 & 14,0 & 2,10\% & 0,10 & 0,16 & 2,95\% \\
Caucho y plástico & 2,5 & 3,1 & 1,45\% & 0,07 & 0,11 & 2,76\% \\
Minerales no metálicos & 9,7 & 10,9 & 0,78\% & 0,08 & 0,10 & 1,39\% \\
Metalurgia y productos metálicos & 5,3 & 6,3 & 1,16\% & 0,17 & 0,21 & 1,38\% \\
Equipo eléctrico y electrónico & 1,7 & 2,9 & 3,72\% & 0,04 & 0,05 & 1,25\% \\
Maquinaria e instalación & 4,1 & 4,0 & -0,16\% & 0,09 & 0,11 & 0,93\% \\
Vehículos y equipo de transporte & 2,1 & 2,2 & 0,48\% & 0,04 & 0,04 & -1,20\% \\
Muebles y colchones & 1,9 & 2,4 & 1,55\% & 0,14 & 0,18 & 1,64\% \\
Otras manufactureras & 3,1 & 3,6 & 1,03\% & 0,09 & 0,09 & -0,48\% \\
Energía eléctrica & 12,7 & 18,5 & 2,56\% & 0,04 & 0,06 & 2,66\% \\
Gas, vapor y aire acondicionado & 2,1 & 3,3 & 2,93\% & 0,01 & 0,02 & 2,96\% \\
Agua & 2,1 & 2,4 & 1,06\% & 0,03 & 0,06 & 4,39\% \\
Aguas residuales, desechos y saneamiento & 3,8 & 4,4 & 0,96\% & 0,03 & 0,06 & 4,26\% \\
Reciclaje & 1,1 & 1,4 & 1,62\% & 0,00 & 0,10 & 27,74\% \\
Edificaciones & 22,0 & 20,5 & -0,49\% & 0,65 & 1,03 & 3,13\% \\
Obras civiles & 9,8 & 11,9 & 1,36\% & 0,12 & 0,21 & 4,05\% \\
Actividades especializadas de construcción & 8,5 & 9,1 & 0,41\% & 0,19 & 0,34 & 3,80\% \\
Comercio; Reparación de vehículos & 51,2 & 89,1 & 3,77\% & 3,14 & 3,90 & 1,45\% \\
Transporte terrestre y tuberías & 24,1 & 31,8 & 1,86\% & 0,88 & 1,30 & 2,66\% \\
Transporte acuático & 0,2 & 0,2 & -0,73\% & 0,01 & 0,02 & 1,84\% \\
Transporte aéreo & 2,9 & 7,7 & 6,75\% & 0,02 & 0,03 & 3,63\% \\
Almacenamiento y apoyo al transporte & 4,9 & 9,5 & 4,50\% & 0,18 & 0,26 & 2,43\% \\
Correo y mensajería & 0,7 & 1,2 & 3,41\% & 0,07 & 0,19 & 7,27\% \\
Alojamiento y comida & 23,8 & 40,6 & 3,63\% & 0,86 & 1,73 & 4,76\% \\
Información y comunicaciones & 18,3 & 31,3 & 3,66\% & 0,40 & 0,40 & -0,07\% \\
Financieras y seguros & 22,3 & 53,1 & 5,95\% & 0,22 & 0,43 & 4,49\% \\
Inmobiliarias & 59,9 & 90,1 & 2,75\% & 0,21 & 0,32 & 2,64\% \\
Profesionales, científicas y técnicas & 21,8 & 31,2 & 2,43\% & 0,30 & 0,57 & 4,43\% \\
Servicios administrativos y de apoyo & 23,6 & 38,8 & 3,38\% & 0,54 & 1,08 & 4,76\% \\
Administración pública y defensa & 31,3 & 62,7 & 4,75\% & 0,50 & 0,81 & 3,30\% \\
Educación de mercado; Educación de no mercado & 31,1 & 52,8 & 3,60\% & 0,75 & 1,01 & 2,01\% \\
Salud y servicios sociales & 23,6 & 49,8 & 5,11\% & 0,70 & 1,01 & 2,51\% \\
Arte, recreación y otros servicios & 11,0 & 40,6 & 9,09\% & 0,97 & 1,24 & 1,65\% \\
Hogares como empleadores & 4,3 & 6,7 & 3,06\% & 0,67 & 0,72 & 0,42\% \\
\bottomrule
\end{longtable}
\endgroup
{\footnotesize Nota: El renglón Alimentos, bebidas y tabaco agrega manufacturas de carnes y pescado, aceites, lácteos, molinería y panadería, café, azúcar y panela, cacao y confitería, frutas y hortalizas, otros alimentos, bebidas y tabaco. Fuente: cálculos propios con DANE y GEIH.}


%\bibliographystyle{apalike}
%\bibliography{Paper/references}

\end{document}
